% Generated draft: requires mathematical and source review.
\documentclass[11pt]{article}
\usepackage[margin=25mm]{geometry}
\usepackage{fontspec}
\setmainfont{Latin Modern Roman}
\usepackage{amsmath,amssymb,mathtools,unicode-math}
\setmathfont{Latin Modern Math}
\usepackage{enumitem,needspace,xurl,hyperref}
\setlength{\parindent}{0pt}
\setlength{\parskip}{5pt}
\setlength{\emergencystretch}{3em}
\newcommand{\N}{\mathbb{N}}
\newcommand{\Z}{\mathbb{Z}}
\newcommand{\Q}{\mathbb{Q}}
\newcommand{\R}{\mathbb{R}}
\newcommand{\C}{\mathbb{C}}
\newcommand{\F}{\mathbb{F}}
\newcommand{\A}{\mathbb{A}}
\newcommand{\PP}{\mathbb{P}}
\newcommand{\E}{\mathbb{E}}
\newcommand{\eps}{\varepsilon}
\newcommand{\dd}{\,\mathrm{d}}
\begin{document}
\section*{UnsolvedMath ID: 800944}
\textbf{Source problem code:} catalogue entry
\subsection{Definitions and mathematical statement}
Let \(M\) be a closed connected surface of genus \(h\geq 2\), and let \(g\) be a Riemannian metric whose Gaussian curvature satisfies
\[
K_g(x)\leq 0\qquad\text{for every }x\in M.
\]
Here closed means compact without boundary. We use the conventional interpretation that the surface is orientable and the metric is smooth; the supplied statement does not explicitly specify orientability or metric regularity, so these conventions require review.

Define the unit tangent bundle by
\[
SM=\{(x,v)\in TM:g_x(v,v)=1\}.
\]
For \((x,v)\in SM\), let \(\gamma_{x,v}\) be the geodesic with initial position \(x\) and initial velocity \(v\). The geodesic flow is
\[
\varphi_t(x,v)=\bigl(\gamma_{x,v}(t),\dot\gamma_{x,v}(t)\bigr),\qquad t\in\R.
\]
Let \(\pi:SM\to M\) be the bundle projection and define the canonical contact form by
\[
\alpha_{(x,v)}(\xi)=g_x\bigl(v,\dd\pi_{(x,v)}\xi\bigr).
\]
The normalized Liouville measure is the flow-invariant probability measure
\[
\mu_L(A)=\frac{\displaystyle\int_A|\alpha\wedge\dd\alpha|}{\displaystyle\int_{SM}|\alpha\wedge\dd\alpha|}
\]
for measurable \(A\subseteq SM\). Normalization does not change ergodicity.

\textbf{Question 24 (Burns--Matveev).} For every surface and metric satisfying these hypotheses, is \(\varphi_t\) ergodic with respect to \(\mu_L\)? Explicitly, must every measurable \(A\subseteq SM\) satisfying
\[
\mu_L\bigl(\varphi_t^{-1}(A)\mathbin{\triangle}A\bigr)=0\qquad\text{for every }t\in\R
\]
satisfy \(\mu_L(A)\in\{0,1\}\), where \(\triangle\) denotes symmetric difference? The curvature hypothesis permits zero curvature; no strict negativity or additional hypothesis on the zero-curvature set is supplied.
\subsection{Short English statement}
For every closed surface of genus at least two with a nonpositively curved Riemannian metric, must its geodesic flow be ergodic for Liouville measure? The formulation uses the conventional smooth, orientable setting, which the supplied wording does not explicitly specify. The local record labels the question open; that label is not independently verified.
\subsection{Sources}
\begin{itemize}
\item Local UnsolvedMath registry, permanent ID \(800944\), source problem code \texttt{catalogue entry}, title \emph{Open Problems and Questions about Geodesics --- Question 24}. The record attributes the question to Burns and Matveev, gives proposed year 2019, and labels its status \texttt{open}.
\item Burns and Matveev, \emph{Open Problems and Questions about Geodesics} (2019, as identified in the supplied record), Question 24. Supplied source URL: \url{https://arxiv.org/abs/1308.5417}. The record reports access on 2026-10-16 and extraction method \texttt{arxiv\_latex}. These are supplied provenance details; the source was not independently visited or verified for this draft.
\end{itemize}
\end{document}
